\begin{abstrak}
Perkembangan aplikasi mobile membuka peluang baru dalam penyediaan tes psikologi yang lebih mudah diakses dan efisien dibandingkan metode konvensional, namun tantangan utama terletak pada bagaimana menerjemahkan desain antarmuka yang kompleks ke dalam aplikasi yang konsisten, responsif, dan mudah digunakan. Aplikasi Jatidiri dikembangkan sebagai media penyajian berbagai tes psikologi, seperti tes kecerdasan, kendali diri, bakat, kebahagiaan, kecemasan, gaya belajar, kendali stres, kesehatan psikologis, potensi, karier, kecerdasan dewasa, dan ketangguhan, yang ditujukan untuk membantu pengguna memahami kondisi psikologis dan potensi diri mereka secara digital. Penelitian ini bertujuan untuk merancang dan mengimplementasikan antarmuka aplikasi Jatidiri berbasis mobile menggunakan framework Flutter berdasarkan desain yang dibuat di Figma, dengan menerapkan metode Prototype agar proses pengembangan dapat dilakukan secara iteratif dan adaptif terhadap umpan balik.

Metode Prototype digunakan melalui tahapan pengumpulan kebutuhan, pembuatan prototype antarmuka, evaluasi dan penyempurnaan prototype, serta implementasi dan pengujian akhir. Fokus utama pengembangan adalah penerjemahan desain UI dari Figma ke dalam komponen Flutter dengan memperhatikan konsistensi tata letak, warna, tipografi, dan interaksi pengguna. Pengujian dilakukan menggunakan metode black box untuk memastikan setiap fitur dan tampilan antarmuka berfungsi sesuai dengan kebutuhan dan desain yang telah ditetapkan.

Hasil penelitian menunjukkan bahwa metode Prototype efektif dalam mendukung proses implementasi desain Figma ke dalam aplikasi Flutter, sehingga menghasilkan antarmuka aplikasi yang responsif, konsisten dengan rancangan awal, dan mudah digunakan. Namun, penelitian ini memiliki keterbatasan karena tidak mencakup validasi psikometrik terhadap isi tes maupun evaluasi usability secara formal terhadap pengguna akhir, sehingga hasil penelitian difokuskan pada aspek perancangan dan implementasi teknis antarmuka aplikasi.
\end{abstrak}

\noindent \textbf{Kata kunci:} Aplikasi Mobile; Tes Psikologi; Flutter; Metode Prototype; UI/UX.
